\chapter{Marco metodológico}
\label{chap:marco_metodologico}

\section{Posicionamiento metodológico: el experimento de enseñanza}

Esta investigación adopta el experimento de enseñanza como diseño metodológico. La elección responde a la pregunta que la orienta: ¿cómo diseñar e implementar tareas de trigonometría con inteligencia artificial de manera que contribuyan al desarrollo del pensamiento trigonométrico en estudiantes de grado décimo?

Esa pregunta no se responde midiendo resultados de aprendizaje en dos grupos paralelos. Se responde observando con detenimiento cómo piensan los estudiantes cuando resuelven esas tareas; qué errores cometen y qué revelan sobre su comprensión; cómo interactúan con la IA y qué papel juega esa interacción en su razonamiento; y de qué manera el diseño de las tareas facilita —o no— el tránsito entre representaciones gráficas, algebraicas y contextuales de las funciones trigonométricas.

El experimento de enseñanza ofrece algo que los diseños comparativos no pueden: acceso al pensamiento matemático \emph{mientras se construye} \parencite{steffe2000teaching}. El investigador no observa desde afuera el resultado de un proceso educativo ya terminado. Participa activamente en él: diseña tareas, las implementa con los estudiantes, registra lo que ocurre, y vuelve al diseño con ese conocimiento acumulado. Esa iteratividad no es una característica secundaria del método. Es su lógica central.

\textcite{gonzalezAcercamientoInvestigacionDiseno2011} precisan que en los experimentos de enseñanza el objeto de estudio no es el estudiante en abstracto sino los procesos de aprendizaje matemático en un contexto específico de interacción: con el docente, con las tareas, con las herramientas disponibles. En esta investigación, esa herramienta es la inteligencia artificial generativa. Documentar cómo media el pensamiento trigonométrico —cuándo abre el pensamiento, cuándo lo reemplaza, en qué condiciones resulta andamiaje genuino— es el núcleo del trabajo.

\textcite{cobb2003design} sitúan los experimentos de enseñanza dentro de una familia más amplia de diseños que comparten un rasgo: el objeto de estudio no es independiente del proceso de investigación, sino que se construye a través de él. Las tareas no se diseñan antes de la investigación y se aplican mecánicamente. Se diseñan, se implementan, se analizan, y se rediseñan. Eso es lo que aquí ocurrirá.

\subsection{Por qué no un diseño comparativo}

Una aproximación razonable sería comparar dos grupos: uno con acceso a IA, otro sin acceso. Esa elección tiene valor propio y responde preguntas del tipo: ¿qué diferencias en desempeño o en estrategias emergen cuando se tiene o no se tiene IA disponible?

Pero eso no es lo que aquí interesa. Lo que interesa es comprender \emph{cómo} —bajo qué condiciones de diseño, con qué tipo de consignas, en qué momentos específicos de la secuencia— la IA puede funcionar como agente de acompañamiento matemático que provoca pensamiento en lugar de reemplazarlo. Esa pregunta exige diseño iterativo, observación sostenida, y análisis que capture la particularidad de los procesos. No exige comparación estadística.

\section{Paradigma de investigación}

Esta investigación se posiciona en el paradigma pospositivista \parencite{phillipsBurbules2000, creswell2014research}. El pospositivismo reconoce que existe una realidad —en este caso, el pensamiento trigonométrico de los estudiantes y los procesos de aprendizaje que ocurren en el aula— pero acepta que esa realidad solo puede conocerse de forma aproximada e imperfecta. A diferencia del positivismo clásico, que asume objetividad plena y medición exacta, el pospositivismo reconoce que el investigador influye sobre lo que estudia y que el conocimiento siempre es provisional. A diferencia del paradigma interpretativo, no renuncia a la búsqueda de patrones sistemáticos ni al rigor metodológico como criterio de validez.

Para esta investigación, ese posicionamiento tiene implicaciones concretas. El pensamiento trigonométrico de los estudiantes existe como objeto real —no es construido arbitrariamente por el investigador—, pero solo puede aproximarse a través de múltiples fuentes de datos (observaciones, entrevistas, producciones escritas, registros de interacción con la IA) analizadas con rigor y trianguladas. Ninguna fuente individual ofrece acceso completo a lo que realmente ocurre en la mente del estudiante cuando resuelve una tarea. Las fuentes combinadas permiten aproximaciones más confiables.

\textcite{creswell2014research} señala que el pospositivista trabaja con métodos cualitativos cuando la naturaleza del fenómeno así lo exige: cuando lo que importa comprender es el proceso y no solo el resultado, cuando los participantes son pocos y el contexto es específico, cuando el análisis estadístico no puede capturar lo que se busca. Todas esas condiciones se cumplen aquí.

\textcite{camargoConferenciaInteramericanaEducacion10} precisa que en educación matemática las estrategias de diseño —como el experimento de enseñanza— no buscan neutralidad del investigador sino rigor en la documentación de sus decisiones y sus efectos. El investigador-docente no se borra del proceso. Lo registra sistemáticamente y lo convierte en parte del análisis.

\section{Fundamentación teórica de las decisiones metodológicas}

Tres marcos teóricos orientan las decisiones de diseño, implementación y análisis de este experimento de enseñanza.

\textbf{Pensamiento trigonométrico.} \textcite{montiel-espinosaDesarrolloPensamientoTrigonometrico2013} desarrolla el concepto de pensamiento trigonométrico como la capacidad de transitar entre la representación geométrica del ángulo, la representación funcional de las razones trigonométricas, y la modelación de fenómenos periódicos reales, sin que ninguno de estos registros sea suficiente por sí solo. Ese tránsito es lo que las tareas de esta investigación buscan provocar y lo que el análisis documentará. La pregunta no es si el estudiante sabe la fórmula de la función seno. Es si puede moverse con soltura entre la situación física, la gráfica y la expresión algebraica, y si puede justificar por qué cada parámetro tiene el valor que tiene.

\textbf{Registros de representación semiótica.} \textcite{duval2006} muestra que la comprensión matemática no reside en dominar un registro aislado —la ecuación algebraica, la gráfica, la tabla numérica— sino en coordinar activamente entre ellos. Los errores más persistentes en trigonometría escolar emergen precisamente cuando los estudiantes no logran articular lo que la gráfica dice con lo que la ecuación representa. Las tareas de esta investigación están diseñadas para hacer visible esa articulación —o su ausencia—, y las observaciones registrarán específicamente cuándo y cómo los estudiantes alternan o fallan en alternar entre registros.

\textbf{Diseño de tareas matemáticas.} \textcite{leungDesigningMathematicsTasks2015} identifican criterios para tareas que exigen pensamiento matemático profundo: deben admitir múltiples caminos de solución, conectar con fenómenos reales, y crear las condiciones para que errores conceptuales emerjan y sean examinados. \textcite{radmehrTheoreticalFrameworkTask2023} y \textcite{canogullariTaskDesignPrinciples2025} añaden que una tarea de alta exigencia cognitiva no puede resolverse aplicando directamente una fórmula memorizada: exige explorar, comparar estrategias y justificar decisiones. Ese criterio rige el diseño de las dos tareas de esta investigación.

La inteligencia artificial entra en este marco no como tema de estudio sino como herramienta de mediación. Su rol no es entregar respuestas correctas. Es funcionar como agente de acompañamiento que el estudiante interroga, evalúa y contrasta con sus propias construcciones. Este criterio —que la IA nunca resuelva directamente sino que provoque pensamiento— es la apuesta pedagógica central del diseño, y documentar cuándo se cumple y cuándo no es uno de los objetivos del análisis.

\section{Diseño del experimento de enseñanza}

\subsection{El docente-investigador}

En el experimento de enseñanza, el investigador no es un observador externo. Es, al mismo tiempo, quien diseña las tareas, las implementa con los estudiantes, registra lo que ocurre, y utiliza ese conocimiento para refinar el diseño. \textcite{steffe2000teaching} denominan a esta figura el \emph{teacher-researcher}: alguien que actúa pedagógicamente y analiza simultáneamente.

Eso tiene implicaciones concretas aquí. Las dos tareas que se describen en este capítulo son la versión actual del diseño. No son definitivas. Si durante la implementación emerge un obstáculo conceptual no anticipado, o si la IA genera un tipo de interacción que modifica el rumbo de la sesión, el diseño se ajustará. Eso no es falta de rigor. Es exactamente lo que el método exige.

\subsection{Contexto e institución}

El experimento se implementa en el Colegio Villa de las Palmas, institución de educación media privada ubicada en Palmira, Valle del Cauca. El municipio cuenta con aproximadamente 300.000 habitantes y se encuentra a 30 km al norte de Cali, en una zona urbana con infraestructura educativa desarrollada.

La institución cuenta con aulas equipadas con computadores e internet, lo que hace técnicamente posible el uso simultáneo de GeoGebra y ChatGPT durante las sesiones de trabajo. Al año 2026, el Colegio Villa de las Palmas tiene una matrícula de 230 estudiantes. La selección de este contexto responde a un criterio práctico y epistemológico a la vez: es la institución donde el investigador ejerce como docente, lo que garantiza acceso sostenido al grupo durante el período de implementación y una familiaridad con el contexto que ningún observador externo podría reemplazar.

\subsection{Participantes}

\textbf{Grupo de trabajo.} Se trabaja con el grupo completo de grado décimo, conformado por 25 estudiantes, cuyas edades oscilan entre 15 y 17 años. Todos los estudiantes del grupo participan en las sesiones de trabajo con las tareas.

\textbf{Muestra intensiva.} Para el análisis más detallado —entrevistas semiestructuradas, análisis de producciones, seguimiento de trayectorias individuales— se seleccionan 4 a 6 estudiantes que representen la diversidad de desempeños del grupo:

\begin{itemize}
  \item Al menos 1-2 estudiantes con desempeño alto en matemáticas.
  \item Al menos 1-2 estudiantes con desempeño medio o variable.
  \item Al menos 1-2 estudiantes con dificultades persistentes en matemáticas.
  \item Representación de ambos géneros en la medida de lo posible.
\end{itemize}

Esta muestra intensiva no es una muestra estadística. Es un mecanismo para documentar en profundidad trayectorias de aprendizaje particulares que ilustren patrones más amplios del grupo.

\subsection{La secuencia de tareas}

Las tareas son el corazón del experimento. Se diseñan dos tareas para la versión inicial de implementación:

\begin{itemize}
  \item \textbf{Tarea 1}: Las transformaciones de la función seno: la rueda de la fortuna.
  \item \textbf{Tarea 2}: Modelar con coseno: el movimiento de las mareas.
\end{itemize}

Ambas tareas responden a cinco criterios didácticos que rigen el diseño \parencite{leungDesigningMathematicsTasks2015, radmehrTheoreticalFrameworkTask2023, canogullariTaskDesignPrinciples2025}:

\begin{enumerate}
  \item \textbf{Exigencia matemática alta}: la tarea no puede resolverse aplicando directamente una fórmula. Exige explorar, comparar estrategias y justificar decisiones.
  \item \textbf{Articulación de representaciones}: cada actividad exige trabajar con al menos dos tipos de representación —gráfica, algebraica, numérica o verbal— y transitar entre ellas \parencite{duval2006}.
  \item \textbf{Rol explícito de la IA}: se define con precisión cuándo y cómo usa el estudiante ChatGPT: para interrogar, contrastar o solicitar retroalimentación; nunca para obtener la respuesta directa de la tarea.
  \item \textbf{Espacio para el error}: las consignas admiten respuestas parciales o erróneas como punto de partida legítimo para la discusión, no como faltas a corregir.
  \item \textbf{Contextualización}: las situaciones conectan con fenómenos periódicos reales —el movimiento de una rueda de la fortuna, las mareas costeras— y tienen sentido fuera del aula de matemáticas.
\end{enumerate}

Cada tarea está organizada en cinco momentos didácticos articulados, cuya secuencia responde a una lógica de construcción progresiva en la que cada momento prepara el siguiente. La Tarea 1 activa el tránsito de la razón trigonométrica a la función seno como objeto matemático propio; la Tarea 2 exige movilizar ese conocimiento para modelar un fenómeno diferente con coseno. El diseño completo de cada tarea, con sus momentos didácticos, errores esperados y criterios de evaluación, se presenta en la sección de métodos y materiales.

\subsection{Ciclos del experimento}

El experimento se organiza en tres ciclos iterativos \parencite{cobb2003design}:

\begin{description}
  \item[Ciclo 1 — Diseño:] Elaboración de las tareas con base en criterios teóricos, revisión de la literatura sobre obstáculos conceptuales en trigonometría, y diseño de instrumentos de recolección de datos (guías de observación, protocolo de entrevistas, rúbrica analítica).
  \item[Ciclo 2 — Implementación y registro:] Aplicación de las tareas con el grupo durante 8 a 10 semanas. Observación sistemática, recolección de producciones escritas, grabaciones de audio o video con consentimiento, y entrevistas con la muestra intensiva.
  \item[Ciclo 3 — Análisis y rediseño:] Análisis del material recolectado, identificación de patrones de razonamiento y obstáculos persistentes, y, si el tiempo y las condiciones lo permiten, rediseño de las tareas con base en los hallazgos para un ciclo posterior.
\end{description}

\section{Recolección de datos}

La recolección se organiza alrededor de documentar cómo los estudiantes abordan las tareas: qué estrategias despliegan, en qué puntos se quedan atrapados, cómo interactúan con la IA, y cómo evoluciona su pensamiento a lo largo de la secuencia.

\textbf{Observación de sesiones.} Durante cada sesión se registra mediante notas de campo lo que ocurre: estrategias que cada estudiante intenta, momentos específicos de bloqueo, preguntas que formula (al docente-investigador o a la IA), representaciones que usa espontáneamente, y manifestaciones de comprensión o confusión. Se documentan especialmente los momentos de consulta a la IA: qué pregunta exacta formula el estudiante, cómo interpreta la respuesta, si cambia de estrategia después. Las sesiones se graban en audio o video cuando sea posible y con consentimiento informado.

\textbf{Entrevistas semiestructuradas.} Con la muestra intensiva (4-6 estudiantes), se realizan entrevistas de 15 a 20 minutos después de que completan conjuntos de tareas. Las conversaciones buscan acceder directamente al pensamiento del estudiante: ¿cuál fue tu primer paso? ¿dónde te atascaste? ¿qué hiciste entonces? ¿cuándo decidiste consultar la IA y por qué? ¿la respuesta que recibiste te resultó útil o te confundió? No se trata de cuestionarios rígidos. El investigador escucha, hace preguntas de seguimiento basadas en lo que el estudiante dice, y evita sugerir respuestas.

\textbf{Producciones de los estudiantes.} Se recolecta y preserva todo lo que escriben y dibujan durante las sesiones: cálculos, gráficos dibujados a mano, tablas de datos, diagramas, borradores con correcciones. Especialmente los borradores: muestran el proceso de pensamiento que la respuesta final borra. También se registran las capturas de las interacciones con ChatGPT. Este material es central para el análisis.

\textbf{Cuestionarios de contexto.} Al inicio se aplica un cuestionario breve (15 minutos) que mide conocimiento previo de trigonometría, autopercepción matemática, y experiencia previa con herramientas de IA. Al final, otro cuestionario de similar extensión mide autopercepción posterior y reflexión sobre la experiencia con la herramienta. Estos cuestionarios no buscan cuantificar aprendizaje. Buscan contexto para interpretar lo que se observó.

\section{Procedimiento de análisis}

El análisis no busca producir números ni verificar hipótesis previas. Busca describir con detalle y rigor cómo piensan realmente estos estudiantes.

\textbf{Fase 1 — Preparación e inmersión.} Se transcribe cada grabación palabra por palabra, se digitalizan fotografías de producciones, y se revisan y amplían las notas de campo. Este acto de organizar ya es análisis: al transcribir emergen patrones; al ampliar notas, se recuperan contextos que los datos brutos no capturan.

\textbf{Fase 2 — Codificación abierta.} Se lee el material sin categorías preconcebidas. La pregunta guía es: ¿qué está ocurriendo realmente aquí? Se marcan episodios relevantes: momentos donde el estudiante muestra comprensión nueva; errores específicos y cómo los intenta corregir; cambios de estrategia; consultas a la IA —qué pregunta exacta, cómo interpreta la respuesta, qué hace después. Las categorías no se fuerzan. Emergen del material.

\textbf{Fase 3 — Codificación axial.} Se agrupan códigos abiertos en categorías mayores. Códigos sobre ``dibuja la situación'', ``construye tabla de valores'' se agrupan bajo ``modos de representación visual''. Se buscan relaciones: ¿existe conexión entre el tipo de error y la estrategia inicial? ¿La IA modificó esa relación?

\textbf{Fase 4 — Análisis de patrones.} Se examinan los patrones emergentes a lo largo de la secuencia completa: ¿cómo evoluciona el pensamiento del Momento 1 al Momento 5? ¿Los obstáculos conceptuales documentados en la literatura —confusión entre amplitud y período, dificultad con el desfase— aparecen con la misma forma que se describe, o se manifiestan de manera diferente en este grupo específico? ¿Cuándo usa el estudiante la IA de forma que expande su razonamiento y cuándo la usa para eludir el pensamiento?

\textbf{Fase 5 — Triangulación.} Se aplican dos tipos de triangulación para fortalecer la credibilidad de los hallazgos \parencite{creswell2014research}:

\textbf{(a) Triangulación de fuentes:} los patrones identificados en las notas de observación se contrastan con las producciones escritas de los estudiantes y con los registros de las entrevistas. Un hallazgo que aparece en las tres fuentes gana credibilidad; una discrepancia entre fuentes —el estudiante afirma que entendió, pero sus cálculos muestran confusión— se examina como dato en sí mismo. Ahí suele estar la comprensión más real.

\textbf{(b) Triangulación de investigador:} los hallazgos preliminares se revisan con los directores de tesis en sesiones de análisis conjunto, reduciendo el riesgo de sesgos interpretativos del investigador-docente. Esta revisión es especialmente necesaria dado que el investigador es también quien implementó las tareas y conoce personalmente a los estudiantes.

\textbf{Fase 6 — Narrativas analíticas.} El producto final no es un listado de códigos ni una tabla de frecuencias. Son narrativas que explican cómo piensan realmente los estudiantes. Con estructura. Con ejemplos específicos tomados del material. Con reconocimiento explícito de lo que no puede concluirse. Ese es el tipo de conocimiento que esta investigación busca producir.

\section{Ubicación del proyecto}
\label{sec:ubicacion_proyecto}

\subsection{Localización geográfica e institucional}

La investigación se desarrolla en el Colegio Villa de las Palmas, institución educativa privada de educación media ubicada en Palmira, Valle del Cauca, Colombia. Palmira es un municipio ubicado en el norte del departamento del Valle del Cauca, aproximadamente a 30 kilómetros de Cali, con una población de cerca de 300,000 habitantes.

El Colegio Villa de las Palmas atiende estudiantes de educación básica secundaria y educación media. Al año 2026 cuenta con una matrícula de 230 estudiantes, distribuidos en los niveles de sexto a undécimo grado. La institución dispone de aulas equipadas con computadores e internet, acceso a recursos tecnológicos, y apertura de la administración para facilitar el desarrollo de investigaciones educativas en sus instalaciones.

\subsection{Población estudiantil}

La población de referencia son estudiantes de grados décimo y undécimo del colegio. Esta población fue seleccionada por varias razones:

\begin{enumerate}
\item \textbf{Madurez cognitiva:} Estudiantes en edades entre 15 y 17 años se encuentran en etapa de pensamiento formal, con capacidades de abstracción y razonamiento lógico-matemático desarrolladas.

\item \textbf{Familiaridad con trigonometría:} Es en décimo y undécimo donde se enseña trigonometría formalmente en el currículo colombiano, por lo que estos estudiantes tienen acceso a instrucción sistemática en el tema.

\item \textbf{Experiencia con tecnología:} Estudiantes de esta edad típicamente tienen experiencia previa con herramientas digitales, lo que reduce la curva de aprendizaje asociada al uso de la herramienta de IA.

\item \textbf{Disponibilidad temporal:} Es posible coordinar con la institución sesiones de trabajo regulares dentro de horario académico.
\end{enumerate}

\subsection{Cronología de la implementación}

La implementación está planeada para una duración de 8 a 10 semanas, distribuidas a lo largo de un período académico semestral. Esto permite:

\begin{itemize}
\item Sesiones de trabajo regulares con estudiantes (idealmente 2-3 sesiones por semana de 45-60 minutos cada una).
\item Tiempo suficiente para que los estudiantes completen un conjunto coherente de tareas trigonométricas.
\item Espacio para entrevistas con estudiantes seleccionados sin interrumpir demasiado el horario de clases.
\item Período de análisis de datos al final sin presión de tiempo extrema.
\end{itemize}

La cronología específica se coordina con la institución considerando calendarios académicos, períodos de evaluación, y disponibilidad de espacios físicos.

\subsection{Justificación de la selección del sitio}

Se seleccionó el Colegio Villa de las Palmas por:

\begin{enumerate}
\item \textbf{Accesibilidad:} La institución está dispuesta a permitir la investigación y facilitar acceso a estudiantes y espacios.

\item \textbf{Características pedagógicas:} Existe disposición del docente de matemáticas a participar, coordinando la implementación sin que esto interfiera significativamente con el currículo regular.

\item \textbf{Composición estudiantil:} La población es heterogénea en términos de desempeño académico en matemáticas, lo que permite documentar cómo estudiantes con distintos niveles de competencia se relacionan con la IA.

\item \textbf{Recursos tecnológicos:} El colegio tiene acceso a computadores y conectividad, lo que permite que el grupo experimental tenga acceso confiable a la herramienta de IA.

\item \textbf{Contexto socioeconómico:} Los estudiantes provienen de contextos socioeconómicos medios, lo que implica familiaridad previa con tecnología pero no ``natividad digital'' extrema que pudiera sesgar los resultados.
\end{enumerate}

\section{Métodos y materiales}
\label{sec:metodos_materiales}

\subsection{Participantes y criterios de selección}

\subsubsection{Grupo de trabajo}

El estudio involucra el grupo completo de grado décimo del Colegio Villa de las Palmas: 25 estudiantes, con edades entre 15 y 17 años. La institución cuenta con 230 estudiantes en total al año 2026.

Los criterios de inclusión son:

\begin{itemize}
\item Estar matriculado en grado décimo en el colegio durante el período de implementación.
\item Contar con presencia durante las sesiones de trabajo (8-10 semanas).
\item Consentimiento informado del estudiante y del acudiente.
\end{itemize}

\subsubsection{Muestra intensiva}

Dentro del grupo, se seleccionan 4 a 6 estudiantes para un seguimiento más profundo. Esta muestra intensiva se elige para representar la diversidad de desempeños presente en el grupo:

\begin{itemize}
\item Al menos 1-2 estudiantes con desempeño alto en matemáticas (calificaciones consistentemente altas en años previos).
\item Al menos 1-2 estudiantes con desempeño medio (calificaciones variables o moderadas).
\item Al menos 1-2 estudiantes con dificultades persistentes en matemáticas.
\item Representación de ambos géneros en la medida de lo posible.
\end{itemize}

Estos 4-6 estudiantes reciben observación especialmente detallada, son entrevistados, y sus producciones se analizan con mayor profundidad. El resto del grupo también es observado, pero con menor intensidad.

El tamaño de la muestra intensiva no responde a criterios estadísticos de representatividad, sino a criterios de riqueza informativa propios de la investigación cualitativa \parencite{creswell2014research}. En coherencia con la lógica del experimento de enseñanza, la suficiencia de la muestra se determina por saturación teórica: el seguimiento continúa hasta que los patrones de razonamiento emergentes dejen de generar categorías analíticas nuevas \parencite{cobb2003design}. La selección de 4 a 6 estudiantes que representen la diversidad de desempeños del grupo garantiza que los patrones identificados reflejen la heterogeneidad real del aula.

\subsection{Diseño de las tareas de trigonometría}

\subsubsection{Propósito y estructura general}

Las tareas son el corazón del experimento de enseñanza. No son ejercicios rutinarios de aplicación de fórmulas. Son problemas diseñados para exigir razonamiento conceptual, toma de decisiones sobre estrategia de solución, y uso flexible de múltiples representaciones \parencite{leungDesigningMathematicsTasks2015, radmehrTheoreticalFrameworkTask2023}.

Se diseñan dos tareas. Cada una:

\begin{enumerate}
\item \textbf{Presenta un problema en contexto:} La tarea no es abstracta. Está incrustada en una situación real que requiere modelación o interpretación:

\begin{itemize}
\item \textbf{Tarea 1 — La rueda de la fortuna:} Una rueda de radio 15 m, con centro a 18 m del suelo, gira completamente cada 40 segundos. Una cabina parte desde el punto más bajo en $t=0$. El estudiante debe construir la función $h(t)$ que describe la altura de la cabina en función del tiempo, justificando cada parámetro.

\item \textbf{Tarea 2 — El movimiento de las mareas:} A partir de datos reales de nivel del mar registrados por una estación costera, el estudiante debe identificar el modelo trigonométrico que describe el fenómeno, determinar período y amplitud en contexto, y usar la función para hacer predicciones verificables.
\end{itemize}

\item \textbf{Requiere movilizar conceptos trigonométricos profundamente \parencite{montiel-espinosaDesarrolloPensamientoTrigonometrico2013}:} La tarea no puede resolverse con una fórmula memorizada. Exige que el estudiante:

\begin{itemize}
\item Interprete qué representa una función trigonométrica en la situación específica.
\item Distinga entre ángulos en grados y radianes.
\item Use y contraste múltiples representaciones: algebraica, gráfica, numérica, verbal \parencite{duval2006}.
\item Razone sobre parámetros —amplitud, período, desplazamiento de fase— y su significado en contexto.
\item Evalúe la validez de su respuesta contrastándola con la situación original.
\end{itemize}

\item \textbf{Admite múltiples caminos de razonamiento:} No existe una única ruta de solución. Un estudiante puede comenzar con una tabla de datos y construir la gráfica; otro puede comenzar algebraicamente y parametrizar; otro puede combinar enfoques. Esto es deliberado: se busca observar la diversidad de estrategias, no forzar conformidad metodológica.

\item \textbf{Está diseñada para revelar dificultades conceptuales:} Cada tarea incorpora obstáculos conceptuales conocidos en la enseñanza de la trigonometría:

\begin{itemize}
\item Confusión entre el valor del ángulo y la medida del lado en un triángulo rectángulo.
\item Interpretación incorrecta de la amplitud en una función periódica.
\item Errores en el cálculo del período o del desplazamiento de fase.
\item Dificultades con radianes frente a grados.
\end{itemize}

Documentar cómo cada estudiante confronta estos obstáculos es uno de los objetivos centrales del análisis.
\end{enumerate}

\subsubsection{Estructura de cada tarea}

Cada tarea se organiza en cinco momentos didácticos articulados, cuyo orden responde a una lógica de construcción progresiva:

\begin{description}
\item[Momento 1 — Exploración inicial (sin IA):] El estudiante observa la función trigonométrica base en GeoGebra y registra sus características. No usa ChatGPT todavía. Lo que construye aquí es la referencia con la que después contrastará lo que la IA le diga.

\item[Momento 2 — Variación de parámetros (sin IA):] El estudiante grafica variaciones de la función ($A\mathrm{sen}(Bx)$, $\mathrm{sen}(x+C)$, etc.) y completa una tabla de comparación. Formula una conjetura sobre qué hace cada parámetro.

\item[Momento 3 — Interrogación a ChatGPT (con IA):] El estudiante consulta a ChatGPT sobre los parámetros que acaba de explorar. Evalúa críticamente lo que la IA responde: ¿coincide con lo que observó? ¿lo contradice? ¿la explicación es completa o parcial?

\item[Momento 4 — Modelación (con o sin IA):] El estudiante construye la función que modela la situación real de la tarea. Usa ChatGPT si lo necesita para explorar o verificar, pero la decisión de qué función construir y por qué es suya.

\item[Momento 5 — Síntesis y validación:] El estudiante justifica por escrito cada parámetro de la función en términos de la situación física. Verifica que la función tenga sentido: $h(0)$ debe coincidir con la condición inicial, el período debe corresponder al dato del problema, etc.
\end{description}

\subsection{La herramienta de inteligencia artificial}

\subsubsection{Herramienta seleccionada}

Se utiliza ChatGPT como herramienta de inteligencia artificial durante la implementación. ChatGPT permite conversación natural en español, tiene conocimiento de trigonometría de nivel secundario, y es accesible en línea desde los computadores del colegio.

\subsubsection{Configuración y rol}

La herramienta se configura mediante instrucciones explícitas al sistema (\emph{system prompt}) que definen el rol pedagógico de la IA durante las tareas. Estas instrucciones son parte integral del diseño del experimento y se aplican antes de cada sesión de trabajo. A continuación se presenta el texto exacto de las instrucciones al sistema:

\bigskip
\begin{quote}
\itshape
Eres un tutor de trigonometría para estudiantes colombianos de grado décimo, con edades entre 15 y 17 años. Tu rol es acompañar el pensamiento matemático del estudiante, no resolverle los problemas.

\textbf{PUEDES:}
\begin{itemize}
\item Responder preguntas conceptuales sobre seno, coseno, tangente, amplitud, período, fase y círculo unitario.
\item Dar retroalimentación sobre los intentos de solución que el estudiante comparte contigo.
\item Formular preguntas que lleven al estudiante a reflexionar sobre lo que hizo y por qué.
\item Sugerir representaciones alternativas: gráfica, tabla de valores, diagrama, lenguaje verbal.
\item Señalar si un razonamiento contiene un error, sin dar la corrección directa.
\item Explicar conceptos de varias maneras cuando el estudiante no comprende la primera explicación.
\end{itemize}

\textbf{NO PUEDES:}
\begin{itemize}
\item Resolver directamente los problemas de las tareas propuestas.
\item Dar la expresión algebraica final de la función que se solicita en la tarea.
\item Realizar todos los cálculos en lugar del estudiante.
\item Confirmar si una respuesta es correcta sin preguntar antes por el razonamiento que la sustenta.
\item Dar la respuesta aunque el estudiante insista en pedirla directamente.
\end{itemize}

Cuando el estudiante pida la respuesta directa, responde con una pregunta que lo dirija a encontrarla por sí mismo. Por ejemplo: ``¿Qué significa el período de una función? ¿Cómo lo calcularías a partir de la gráfica que tienes?''

Usa un lenguaje claro, accesible y respetuoso. Nunca hagas sentir al estudiante que su pregunta es equivocada o trivial.
\end{itemize}
\end{quote}
\bigskip

Este criterio es fundamental. La IA funciona como agente de acompañamiento, no como solucionador. Los estudiantes reciben una instrucción breve sobre cómo usarla: cómo formular preguntas claras, cómo interpretar las respuestas, y en qué momentos tiene sentido consultarla. Se enfatiza que es un recurso disponible, no una obligación.

\subsubsection{Registro de interacciones}

Se registra cada interacción entre estudiante y ChatGPT:

\begin{itemize}
\item La pregunta exacta que formula el estudiante.
\item La respuesta que proporciona la IA.
\item El momento dentro de la tarea en que ocurre la consulta.
\item Las decisiones del estudiante después de la consulta: ¿integró la sugerencia en su trabajo? ¿la ignoró? ¿pidió aclaración?
\end{itemize}

Este registro es fuente de datos central para el análisis. Permite documentar no solo qué sugirió la IA, sino cómo el estudiante procesó esa información.

\subsection{Materiales de trabajo}

Los estudiantes trabajan con los siguientes recursos durante las sesiones:

\begin{itemize}
\item GeoGebra (versión en línea, gratuita) para graficación y exploración dinámica.
\item ChatGPT como agente de acompañamiento matemático en los momentos definidos de cada tarea.
\item Cuaderno o papel para cálculos, borradores y representaciones a mano.
\item Guías impresas de cada tarea (situación, consignas por momento, espacio para registros).
\item Acceso al docente-investigador para preguntas, con la misma restricción que rige para ChatGPT: disponible para orientar, no para resolver.
\end{itemize}

\subsection{Instrumentos de recolección de datos}

\subsubsection{Observación sistemática}

\textbf{Notas de campo.} Durante cada sesión se toman notas detalladas sobre lo que ocurre con cada estudiante de la muestra intensiva, y observaciones generales del grupo:

\begin{itemize}
\item Estrategia inicial que intenta.
\item Representaciones que usa espontáneamente (algebraica, gráfica, tabla, dibujo).
\item Momentos específicos donde se queda bloqueado.
\item Preguntas que formula (al docente o a ChatGPT).
\item Cómo responde a retroalimentación o sugerencias.
\item Cuándo decide consultar la IA, qué tipo de pregunta formula, y qué hace con la respuesta.
\item Cambios de estrategia durante la resolución.
\item Manifestaciones de comprensión o confusión.
\end{itemize}

Las notas se escriben en tiempo real o inmediatamente después de la sesión, mientras los detalles están frescos.

\textbf{Grabación de audio/video.} Cuando sea posible y con consentimiento informado, se graban las sesiones. Esto permite revisar conversaciones con precisión posterior, captar detalles de razonamiento verbal que las notas pueden perder, y disponer de evidencia confiable para los reportes.

\subsubsection{Entrevistas semiestructuradas}

Se realizan entrevistas con la muestra intensiva (4-6 estudiantes) después de completar cada tarea o conjunto de tareas.

\textbf{Duración:} 15 a 20 minutos.

\textbf{Preguntas orientadoras:}

\begin{itemize}
\item ``Cuando empezaste [tarea], ¿cuál fue tu primer paso? ¿Por qué empezaste así?''
\item ``¿En qué momento te sentiste confundido o bloqueado? ¿Qué hiciste entonces?''
\item ``¿Cuál concepto de trigonometría te pareció más difícil en esta tarea?''
\item ``¿Cuándo decidiste consultar ChatGPT? ¿Qué te hizo pensar que lo necesitabas?''
\item ``¿La respuesta que recibiste de la IA fue útil? ¿Por qué? ¿Hiciste algo diferente a lo que sugirió?''
\item ``¿Hay momentos donde ChatGPT te confundió o te llevó por mal camino?''
\item ``¿Hay algo que ahora entiendes sobre trigonometría que no entendías antes de esta tarea?''
\end{itemize}

El investigador escucha activamente, hace preguntas de seguimiento basadas en lo que el estudiante dice, y evita sugerir respuestas.

\subsubsection{Producciones de los estudiantes}

Se recolectan y preservan todas las producciones escritas y visuales generadas durante las sesiones:

\begin{itemize}
\item Hojas de trabajo con cálculos, gráficos dibujados a mano, tablas de datos.
\item Borradores y trabajo sucio —no solo las respuestas finales. Los borradores muestran el proceso de pensamiento que la respuesta limpia borra.
\item Correcciones: cómo y dónde el estudiante modificó su trabajo.
\item Explicaciones escritas del razonamiento.
\item Capturas de pantalla o transcripciones de las interacciones con ChatGPT.
\end{itemize}

El análisis de estas producciones permite ver las estrategias reales usadas (no lo que el estudiante dice que hizo), los errores específicos cometidos y sus patrones, las preferencias de representación, y la evolución del trabajo de la Tarea 1 a la Tarea 2.

\subsubsection{Cuestionarios de contexto}

\textbf{Pre-test (inicio de la implementación):} Un cuestionario breve (máximo 15 minutos) que mide:

\begin{itemize}
\item Conocimiento previo de trigonometría: razones trigonométricas, funciones trigonométricas, período, amplitud, radianes.
\item Autopercepción matemática: ``¿Cuánta confianza tienes en tu capacidad de resolver problemas de trigonometría?'' (escala 1-10).
\item Experiencia previa con herramientas de IA: ``¿Has usado inteligencia artificial antes? ¿Para qué?''
\end{itemize}

\textbf{Post-test (al finalizar la implementación):} Un cuestionario de extensión similar que mide:

\begin{itemize}
\item Conocimiento posterior: preguntas análogas al pre-test sobre los conceptos tratados en las tareas.
\item Cambio en autopercepción: ``¿Cómo cambió tu confianza en trigonometría durante este proceso?''
\item Reflexión sobre la herramienta de IA: ``¿Qué tan útil encontraste ChatGPT? ¿En qué momentos fue más valioso? ¿En qué momentos fue confuso o no te ayudó?''
\item Reflexión sobre las tareas: ``¿Cuál fue el momento más difícil de resolver? ¿Qué recursos usaste?''
\end{itemize}

Los cuestionarios no buscan cuantificar aprendizaje. Proporcionan contexto para interpretar lo que se observó.

\subsection{Procedimientos de implementación}

\textbf{Fase 0 (Semana previa):} Preparación logística. Coordinación con el colegio, obtención de consentimientos informados, verificación de acceso a GeoGebra y ChatGPT desde los computadores del aula, y ajuste final del diseño de las tareas.

\textbf{Fase 1 (Semana 1):} Introducción y línea base. Aplicación del pre-test. Sesión introductoria donde se explica el proyecto a los estudiantes. Familiarización breve con GeoGebra y con el uso de ChatGPT como agente de acompañamiento matemático. Primeras observaciones.

\textbf{Fase 2 (Semanas 2-6):} Implementación principal. Aplicación de las dos tareas en sesiones de 40 a 60 minutos. Observación sistemática del grupo completo, con seguimiento más detallado de la muestra intensiva. Recolección continua de producciones y registro de interacciones con la IA.

\textbf{Fase 3 (Semana 7-8):} Entrevistas con la muestra intensiva. Post-test. Recolección de materiales pendientes.

\textbf{Fase 4 (Semana 9-10):} Organización del corpus de datos. Transcripción de grabaciones. Digitalización de producciones. Preparación del material para el análisis sistemático.

\section{Cronograma de actividades}
\label{sec:cronograma}

El experimento de enseñanza se organiza en tres fases principales: diseño, implementación y análisis. La implementación en aula abarca 10 semanas; el análisis de datos y la escritura de resultados se desarrollan en las semanas siguientes. El cronograma a continuación refleja la planificación de actividades para el primer año de la maestría.

\begin{table}[htbp]
\centering
\caption{Cronograma de actividades del experimento de enseñanza}
\label{tab:cronograma}
\small
\begin{tabular}{|p{5.5cm}|c|c|c|c|c|c|c|c|c|c|}
\hline
\textbf{Actividad} & \textbf{S1} & \textbf{S2} & \textbf{S3} & \textbf{S4} & \textbf{S5} & \textbf{S6} & \textbf{S7} & \textbf{S8} & \textbf{S9} & \textbf{S10} \\
\hline
\multicolumn{11}{|l|}{\textit{Fase 1: Diseño}} \\
\hline
Revisión de literatura y ajuste del marco teórico & $\bullet$ & $\bullet$ & & & & & & & & \\
\hline
Diseño definitivo de tareas 1 y 2 & $\bullet$ & $\bullet$ & & & & & & & & \\
\hline
Construcción de instrumentos (guía obs., protocolo entrevistas, rúbrica) & $\bullet$ & $\bullet$ & & & & & & & & \\
\hline
Trámite de consentimientos informados & $\bullet$ & $\bullet$ & & & & & & & & \\
\hline
Configuración y prueba de herramientas (GeoGebra + ChatGPT) & & $\bullet$ & $\bullet$ & & & & & & & \\
\hline
\multicolumn{11}{|l|}{\textit{Fase 2: Implementación}} \\
\hline
Pre-test y sesión introductoria & & & $\bullet$ & & & & & & & \\
\hline
Tarea 1 — La rueda de la fortuna (2 sesiones, 40 min c/u) & & & & $\bullet$ & $\bullet$ & & & & & \\
\hline
Primeras entrevistas con muestra intensiva & & & & & $\bullet$ & & & & & \\
\hline
Tarea 2 — El movimiento de las mareas (2 sesiones, 40 min c/u) & & & & & & $\bullet$ & $\bullet$ & & & \\
\hline
Segundas entrevistas con muestra intensiva & & & & & & & $\bullet$ & & & \\
\hline
Post-test y entrevistas finales & & & & & & & & $\bullet$ & & \\
\hline
\multicolumn{11}{|l|}{\textit{Fase 3: Análisis}} \\
\hline
Organización del corpus: transcripciones, digitalización & & & & & & & & $\bullet$ & $\bullet$ & \\
\hline
Codificación abierta y axial & & & & & & & & & $\bullet$ & $\bullet$ \\
\hline
Análisis de patrones y triangulación & & & & & & & & & & $\bullet$ \\
\hline
\end{tabular}
\end{table}

\vspace{0.5em}

Las semanas 1--2 corresponden a la fase de preparación previa al inicio de la implementación en aula. Las semanas 3--8 constituyen el período de implementación activa. Las semanas 8--10 inician el análisis, que se extiende en el segundo semestre del año hasta la producción de los resultados y la escritura del informe final.

Este cronograma es indicativo. El experimento de enseñanza, por su naturaleza iterativa, admite ajustes en la secuencia cuando los datos recolectados lo justifiquen.

\section{Presupuesto}
\label{sec:presupuesto}

El presupuesto del proyecto se organiza en cuatro categorías: personal, materiales, servicios digitales y difusión. Dado el carácter cualitativo del estudio y su implementación en una institución donde el investigador ejerce como docente, los costos directos son moderados.

\begin{table}[htbp]
\centering
\caption{Presupuesto estimado del proyecto de investigación}
\label{tab:presupuesto}
\small
\begin{tabular}{|p{5.5cm}|p{6cm}|}
\hline
\textbf{Rubro} & \textbf{Descripción} \\
\hline
\multicolumn{2}{|l|}{\textit{Personal}} \\
\hline
Investigador principal & Tiempo de diseño, trabajo de campo, análisis y escritura. Cubierto como parte de las actividades de la maestría. \\
\hline
\multicolumn{2}{|l|}{\textit{Materiales}} \\
\hline
Material impreso para estudiantes & Guías de tareas, cuestionarios y rúbricas (2 tareas, aproximadamente 20 estudiantes). \\
\hline
Papelería general & Papel, bolígrafos, carpetas para archivo de producciones estudiantiles. \\
\hline
Almacenamiento digital & Dispositivo para copia de seguridad de grabaciones y producciones digitalizadas. \\
\hline
\multicolumn{2}{|l|}{\textit{Servicios digitales}} \\
\hline
ChatGPT & Acceso a la herramienta de IA durante el período de implementación. \\
\hline
GeoGebra & Versión gratuita disponible en línea. Sin costo. \\
\hline
Dispositivo de grabación & Registro de audio de sesiones. Equipo propio del investigador. \\
\hline
\multicolumn{2}{|l|}{\textit{Transporte}} \\
\hline
Desplazamiento al colegio & Transporte durante las sesiones de implementación (10 semanas). \\
\hline
\multicolumn{2}{|l|}{\textit{Difusión}} \\
\hline
Impresión del informe final & Ejemplares para directoras de tesis y depósito en repositorio UNAL. \\
\hline
\end{tabular}
\end{table}

\textbf{Fuente de financiamiento.} El proyecto se financia con recursos propios del investigador. No se solicitan recursos a la Universidad Nacional de Colombia para esta etapa de propuesta. La posibilidad de acceder a convocatorias internas de financiación de la sede Palmira se evaluará al avanzar hacia la fase de implementación.

\section{Consideraciones éticas}
\label{sec:consideraciones_eticas}

La investigación involucra como participantes a estudiantes menores de edad (15 a 17 años), lo que exige el cumplimiento riguroso de procedimientos éticos orientados a proteger su integridad, su privacidad y su participación libre e informada. Los procedimientos descritos a continuación se aplican durante todas las fases del estudio.

\subsection{Consentimiento y asentimiento informado}

Antes del inicio de cualquier actividad de recolección de datos se gestionan tres documentos de consentimiento:

\begin{itemize}
\item \textbf{Consentimiento de padres o acudientes:} documento escrito dirigido al acudiente de cada estudiante, que explica los objetivos del estudio, las actividades en las que participará el estudiante, el tipo de datos que se recolectarán (observaciones, producciones escritas, grabaciones de audio o video, entrevistas), y el uso exclusivamente investigativo de esa información. La firma del documento es condición necesaria para que el estudiante participe en las actividades registradas.

\item \textbf{Asentimiento del estudiante:} documento escrito dirigido al propio estudiante, redactado en lenguaje accesible, que le explica en qué consiste su participación, que puede retirarse en cualquier momento sin consecuencias para su desempeño académico, y que tiene derecho a no responder preguntas que no desee responder en las entrevistas.

\item \textbf{Autorización institucional:} comunicación formal dirigida al rector o coordinador académico del Colegio Villa de las Palmas, que informa sobre el propósito de la investigación, el período de implementación, y el compromiso de compartir los resultados con la institución al finalizar el estudio.
\end{itemize}

\subsection{Confidencialidad y anonimato}

La identidad de los estudiantes no aparece en ningún reporte, análisis o publicación derivada de esta investigación. Para proteger el anonimato:

\begin{itemize}
\item Cada estudiante es identificado mediante un código alfanumérico (E01, E02\ldots E25) que se asigna al inicio de la implementación y se usa de forma consistente en todos los documentos de análisis.
\item Los registros de audio, video y producciones escritas se almacenan en dispositivos de acceso restringido, protegidos con contraseña, a los que solo tiene acceso el investigador principal.
\item Los nombres de los estudiantes no aparecen en transcripciones, notas de campo ni materiales de análisis.
\end{itemize}

\subsection{Manejo de grabaciones}

Las sesiones de trabajo podrán ser grabadas en audio o video únicamente cuando exista consentimiento firmado del acudiente y asentimiento del estudiante. Las grabaciones:

\begin{itemize}
\item Se utilizan exclusivamente para los fines de esta investigación (transcripción, análisis de procesos de razonamiento).
\item No se publican ni comparten en plataformas digitales.
\item Serán eliminadas de todos los dispositivos de almacenamiento en un plazo no mayor a dos años después de la defensa del trabajo de grado.
\end{itemize}

\subsection{Participación voluntaria y sin consecuencias académicas}

La participación en el estudio es voluntaria. Un estudiante puede:

\begin{itemize}
\item Decidir no ser grabado sin que ello afecte su participación en las actividades de clase.
\item Retirarse de las entrevistas o de cualquier actividad de recolección de datos en el momento que lo decida.
\item Solicitar que sus producciones escritas no sean incluidas en el análisis.
\end{itemize}

En ningún caso la decisión de no participar en actividades de recolección de datos afectará la calificación del estudiante en la asignatura de matemáticas.

\subsection{Riesgo mínimo}

Esta investigación se clasifica como de \textbf{riesgo mínimo}, de acuerdo con los criterios establecidos para investigación con seres humanos en Colombia. Las actividades que se desarrollan (resolución de tareas matemáticas, uso de software educativo, conversación con un sistema de IA en contexto académico) no representan riesgo para la integridad física, psicológica o social de los participantes. No se realiza ninguna intervención clínica, diagnóstica ni invasiva. El investigador toma las precauciones necesarias para que el ambiente de trabajo sea seguro, respetuoso y libre de presión.

